% ==================================================================
%  徐森林、金亚东、薛春华《数学分析》（第三册）
%  S14.2  函数的幂级数展开式
%  源：原书印刷页 146-161（PDF 物理页 152-167）
%  底本：PaddleOCR-VL-1.6，见 chapters/_ch14_14_2_raw.md
%  约定：全西文标点；行间公式用一对方括号；多行公式用 align/aligned；
%        跨册引用一律 \extref；数学模式内的中文必须包 \text{}
%  已订正的 OCR 错误与原书照录清单见 tools/_make_ch14_2a.py 头部说明
% ==================================================================
\section{函数的幂级数展开式}\label{sec:14.2}

14.1 节讨论了由幂级数所确定的和函数的性质. 这一节主要讨论一个函数在什么条件下能展开成收敛的幂级数以及如何展开.

\begin{theorem}[幂级数展开的惟一性]\label{thm:14.2.1}
设 $f(x)$ 在 $(x_{0}-R,x_{0}+R)$ 中能展开为收敛的幂级数
\[
  f(x)=\sum_{n=0}^{\infty}a_{n}(x-x_{0})^{n},
\]
则 $f(x)$ 在 $(x_{0}-R,x_{0}+R)$ 中为 $C^{\infty}$ 函数, 且
\[
  a_{n}=\frac{f^{(n)}(x_{0})}{n!},\quad n=0,1,2,\cdots,
\]
即
\[
  f(x)=\sum_{n=0}^{\infty}\frac{f^{(n)}(x_{0})}{n!}(x-x_{0})^{n},
\]
其中 $f^{(0)}(x)=f(x)$.
\end{theorem}

\begin{proof}
由定理~\ref{thm:14.1.4}(3) 知, $f(x)$ 在 $(x_{0}-R,x_{0}+R)$ 中为 $C^{\infty}$ 函数. 又因为
\[
  f^{(k)}(x)=\sum_{n=k}^{\infty}n(n-1)\cdots(n-k+1)a_{n}(x-x_{0})^{n-k},
\]
令 $x=x_{0}$, 得
\[
  f^{(k)}(x_{0})=k(k-1)\cdots2\cdot1\cdot a_{k}=k!a_{k},
\]
所以
\[
  a_{k}=\frac{f^{(k)}(x_{0})}{k!},\quad k=0,1,2,\cdots,
\]
\[
  f(x)=\sum_{n=0}^{\infty}\frac{f^{(n)}(x_{0})}{n!}(x-x_{0})^{n},\quad x\in(x_{0}-R,x_{0}+R).
\]
\end{proof}

\begin{definition}\label{def:14.2.1}
设 $f(x)$ 在 $x=x_{0}$ 处有各阶导数, 于是从 $f(x)$ 可作出一个形式幂级数
\[
  \sum_{n=0}^{\infty}\frac{f^{(n)}(x_{0})}{n!}(x-x_{0})^{n},
\]
称它为 $f(x)$ 在 $x=x_{0}$ 处的 Taylor 级数, 记为
\[
  f(x)\sim\sum_{n=0}^{\infty}\frac{f^{(n)}(x_{0})}{n!}(x-x_{0})^{n}.
\]
\end{definition}

特别当 $x_{0}=0$ 时, 幂级数
\[
  \sum_{n=0}^{\infty}\frac{f^{(n)}(0)}{n!}x^{n}
\]
也称为 $f(x)$ 的 Maclaurin 级数.

如果 $f(x)=\sum_{n=0}^{\infty}\frac{f^{(n)}(x_{0})}{n!}(x-x_{0})^{n}$, $x\in(x_{0}-R,x_{0}+R)$, 则称 $f(x)$ 在 $(x_{0}-R,x_{0}+R)$ 可 Taylor 展开.

\begin{definition}\label{def:14.2.2}
设 $x_{0}\in(a,b)$, 如果 $\exists\delta>0$, 使得幂级数
\[
  \sum_{n=0}^{\infty}a_{n}(x-x_{0})^{n}
\]
在 $(x_{0}-\delta,x_{0}+\delta)$ 中收敛于函数 $f(x)$, 则称 $f(x)$ 在点 $x_{0}$ 处是实解析的; 如果 $\forall x_{0}\in(a,b)$, $f(x)$ 在 $x_{0}$ 处都实解析, 则称 $f(x)$ 在 $(a,b)$ 上为实解析函数或 $f(x)$ 在 $(a,b)$ 上是 $C^{\omega}$ 的函数, 记从 $(a,b)$ 到 $\mathbb{R}$ 的实解析 ($C^{\omega}$) 函数全体为 $C^{\omega}((a,b),\mathbb{R})$, 简记为 $C^{\omega}((a,b))$. $f(x)\in C^{\omega}((a,b),\mathbb{R})$ 或 $f(x)\in C^{\omega}((a,b))$, 或更简单 $f(x)\in C^{\omega}$ 都表示 $f(x)$ 为实解析 ($C^{\omega}$) 函数.
\end{definition}

\begin{remark}\label{rem:14.2.1}
只要 $f(x)$ 在 $x=x_{0}$ 处有任意阶导数, 就能作出它的 Taylor 级数
\[
  \sum_{n=0}^{\infty}\frac{f^{(n)}(x_{0})}{n!}(x-x_{0})^{n},
\]
但这个级数不一定是收敛的 (例~\ref{ex:14.2.1}), 它只是一个形式的幂级数, 不必在意它! 即使收敛, 也未必收敛到 $f(x)$ 自己 (例~\ref{ex:14.2.2}), 对这种 Taylor 级数也不必太在意它!
\end{remark}

\begin{example}\label{ex:14.2.1}
设
\[
  f(x)=\sum_{n=0}^{\infty}\frac{\sin2^{n}x}{n!}
\]
为由函数项级数确定的函数, 因为
\[
  u_{n}^{(k)}(x)=\left(\frac{\sin2^{n}x}{n!}\right)^{(k)}\xlongequal{\text{归纳}}\frac{(2^{n})^{k}\sin\left(2^{n}x+\frac{k\pi}{2}\right)}{n!},
\]
\[
  |u_{n}^{(k)}(x)|\leqslant\frac{(2^{n})^{k}}{n!}=\frac{(2^{k})^{n}}{n!},
\]
$\frac{(2^{k})^{n+1}}{(n+1)!}\bigg/\frac{(2^{k})^{n}}{n!}=\frac{2^{k}}{n+1}\to0$ ($n\to+\infty$), 所以根据正项级数的 $d'$Alembert 判别法的极限形式知, $\sum_{n=0}^{\infty}\frac{(2^{k})^{n}}{n!}$ 收敛 (于 $\mathrm{e}^{2^{k}}$). 再根据函数项级数的 Weierstrass 判别法得到
\[
  \sum_{n=0}^{\infty}u_{n}^{(k)}(x)
\]
在 $(-\infty,+\infty)$ 上一致收敛, 从而推得
\[
  f^{(k)}(x)=\left(\sum_{n=0}^{\infty}u_{n}(x)\right)^{(k)}=\sum_{n=0}^{\infty}u_{n}^{(k)}(x)=\sum_{n=0}^{\infty}\frac{(2^{n})^{k}\sin\left(2^{n}x+\frac{k\pi}{2}\right)}{n!}.
\]
也就是说, $f(x)$ 在 $(-\infty,+\infty)$ 上有任意阶导数, 即 $f(x)$ 为 $C^{\infty}$ 函数. 将 $x=0$ 代入上式得到
\[
  f^{(k)}(0)=\sum_{n=0}^{\infty}\frac{(2^{n})^{k}\sin\frac{k\pi}{2}}{n!}=\begin{cases}
    \displaystyle(-1)^{l}\mathrm{e}^{2^{2l+1}}, & k=2l+1,\\
    0, & k=2l,
\end{cases}
\]
\[
  f(x)\sim\sum_{n=0}^{\infty}\frac{f^{(k)}(0)}{k!}x^{k}=\sum_{l=0}^{\infty}\frac{(-1)^{l}\mathrm{e}^{2^{l+1}}}{(2l+1)!}x^{2l+1}
\]
($f(x)$ 在 $x_{0}=0$ 处的 Taylor 级数).

由于
\[
  \frac{\mathrm{e}^{2^{2l+3}}|x|^{2l+3}}{(2l+3)!}\bigg/\frac{\mathrm{e}^{2^{2l+1}}|x|^{2l+1}}{(2l+1)!}=\frac{\mathrm{e}^{3\cdot2^{2l+1}}x^{2}}{(2l+2)(2l+3)}\to\begin{cases}
    +\infty, & x\neq0,\\
    0, & x=0,
\end{cases}
\]
这个函数项级数除 $x=0$ 外处处发散.
\end{example}

\begin{example}\label{ex:14.2.2}
在\extref{例 3.2.7} 中已证
\[
  f(x)=\begin{cases}
    \mathrm{e}^{-\frac{1}{x}}, & x>0,\\
    0, & x\leqslant0
\end{cases}
\]
为 $C^{\infty}$ 函数, 且 $f^{(n)}(0)=0$, $n=0,1,2,\cdots$. 显然, $f(x)$ 的 Taylor 级数
\[
  \sum_{n=0}^{\infty}\frac{f^{(n)}(0)}{n!}x^{n}=\sum_{n=0}^{\infty}\frac{0}{n!}x^{n}
\]
收敛于 $0$, 而不收敛于 $f(x)$ 自己.
\end{example}

\begin{theorem}\label{thm:14.2.2}
$(a,b)$ 上的 $C^{\omega}$ (实解析) 函数必为 $C^{\infty}$ 函数, 但反之不真.
\end{theorem}

\begin{proof}
设 $f(x)$ 为 $(a,b)$ 上为 $C^{\omega}$ 函数, 则 $\forall x_{0}\in(a,b)$, $\exists\delta>0$, 使得
\[
  \begin{aligned}
    f(x)&=\sum_{n=0}^{\infty}a_{n}(x-x_{0})^{n}\\
    &\xlongequal{\text{定理~\ref{thm:14.2.1}}}\sum_{n=0}^{\infty}\frac{f^{(n)}(x_{0})}{n!}(x-x_{0})^{n},\quad x\in(x_{0}-\delta,x_{0}+\delta)\subset(a,b),
  \end{aligned}
\]
根据定理~\ref{thm:13.2.4p}, $f(x)$ 在 $(x_{0}-\delta,x_{0}+\delta)$ 中是 $C^{\infty}$ 的. 由于 $x_{0}\in(a,b)$ 是任取的, 故 $f(x)$ 为 $(a,b)$ 上的 $C^{\infty}$ 函数.

但反之不真, 例如: \extref{例 3.2.7} 中的
\[
  f(x)=\begin{cases}
    \mathrm{e}^{-\frac{1}{x}}, & x>0,\\
    0, & x<0
\end{cases}
\]
是 $C^{\infty}$ 函数但非 $C^{\omega}$ 函数.
\end{proof}

由定理~\ref{thm:14.2.2} 的证明可看出下面的结论.

\begin{corollary}\label{cor:14.2.1}
可 Taylor 展开的函数或收敛的幂级数函数是 $C^{\infty}$ 的.
\end{corollary}

\begin{theorem}\label{thm:14.2.3}
设 $f(x)$ 在 $(a-R,a+R)$ 中可展开为收敛的幂级数, 即
\[
  f(x)=\sum_{n=0}^{\infty}a_{n}(x-a)^{n},\quad x\in(a-R,a+R).
\]
则 $\forall x_{0}\in(a-R,a+R)$, $f(x)$ 在 $(x_{0}-r,x_{0}+r)$ 中可展开为收敛的幂级数
\[
  f(x)=\sum_{n=0}^{\infty}\frac{f^{(n)}(x_{0})}{n!}(x-x_{0})^{n},
\]
其中 $r=R-|x_{0}-a|$. 于是, $f\in C^{\omega}((a-R,a+R),\mathbb{R})$.
\end{theorem}

\begin{proof}
$\forall x\in(x_{0}-r,x_{0}+r)$, 则
\[
  |x-a|\leqslant|x-x_{0}|+|x_{0}-a|<r+|x_{0}-a|=\left(R-|x_{0}-a|\right)+|x_{0}-a|=R,
\]
所以, 根据定理~\ref{thm:14.1.3}, 有
\[
  \sum_{n=0}^{\infty}\frac{f^{(n)}(a)}{n!}[|x-x_{0}|+|x_{0}-a|]^{n}
\]
绝对收敛, 即
\[
  \begin{aligned}
    \sum_{n=0}^{\infty}\left|\frac{f^{(n)}(a)}{n!}\right|[|x-x_{0}|+|x_{0}-a|]^{n}
    &=\sum_{n=0}^{\infty}\left|\frac{f^{(n)}(a)}{n!}\right|\sum_{k=0}^{n}\mathrm{C}_{n}^{k}|x-x_{0}|^{k}|x_{0}-a|^{n-k}\\
    &=\sum_{k=0}^{\infty}\left[\sum_{n\geqslant k}\left|\frac{f^{(n)}(a)}{n!}\right|\mathrm{C}_{n}^{k}|x_{0}-a|^{n-k}\right]|x-x_{0}|^{k}
  \end{aligned}
\]
收敛. 根据 Cauchy 收敛原理, 有
\[
  \begin{aligned}
    f(x)&=\sum_{n=0}^{\infty}\frac{f^{(n)}(a)}{n!}(x-a)^{n}=\sum_{n=0}^{\infty}\frac{f^{(n)}(a)}{n!}[(x-x_{0})+(x_{0}-a)]^{n}\\
    &=\sum_{n=0}^{\infty}\frac{f^{(n)}(a)}{n!}\sum_{k=0}^{n}\mathrm{C}_{n}^{k}(x-x_{0})^{k}(x_{0}-a)^{n-k}\\
    &=\sum_{k=0}^{\infty}\left[\sum_{n\geqslant k}\frac{f^{(n)}(a)}{n!}\mathrm{C}_{n}^{k}(x_{0}-a)^{n-k}\right](x-x_{0})^{k}\quad(\text{绝对收敛})\\
    &=\sum_{k=0}^{\infty}\frac{1}{k!}\left[\sum_{n\geqslant k}\frac{(f^{(k)})^{n-k}(a)}{(n-k)!}(x_{0}-a)^{n-k}\right](x-x_{0})^{k}\\
    &=\sum_{k=0}^{\infty}\frac{f^{(k)}(x_{0})}{k!}(x-x_{0})^{k},
  \end{aligned}
\]
即 $f(x)$ 在 $(x_{0}-r,x_{0}+r)$ 中可展开为收敛的幂级数. 这就证明了 $f(x)\in C^{\omega}((a-R,a+R),\mathbb{R})$.
\end{proof}

例~\ref{ex:14.2.1} 与例~\ref{ex:14.2.2} 使我们自然产生这样的问题: $f(x)$ 要满足什么条件, 才能保证它的 Taylor 级数收敛于 $f(x)$ 自己? 即在什么条件下, 等式
\[
  f(x)=\sum_{n=0}^{\infty}\frac{f^{(n)}(x_{0})}{n!}(x-x_{0})^{n},\quad x\in(x_{0}-R,x_{0}+R)
\]
成立.

\begin{theorem}[可 Taylor 展开的充分必要条件]\label{thm:14.2.4}
设 $f(x)$ 在 $(x_{0}-R,x_{0}+R)$ 上为 $C^{\infty}$ 函数, 则 $f(x)$ 在 $(x_{0}-R,x_{0}+R)$ 上可 Taylor 展开等价于
\[
  f(x)=\sum_{k=0}^{n}\frac{f^{(k)}(x_{0})}{k!}(x-x_{0})^{k}+R_{n}(x),\quad x\in(x_{0}-R,x_{0}+R),
\]
有
\[
  \lim_{n\to+\infty}R_{n}(x)=0,\quad x\in(x_{0}-R,x_{0}+R).
\]
\end{theorem}

\begin{proof}
$f(x)$ 在 $(x_{0}-R,x_{0}+R)$ 上可 Taylor 展开
\[
  \Leftrightarrow f(x)=\sum_{k=0}^{\infty}\frac{f^{(k)}(x_{0})}{k!}(x-x_{0})^{k}=\lim_{n\rightarrow+\infty}\sum_{k=0}^{n}\frac{f^{(k)}(x_{0})}{k!}(x-x_{0})^{k}
\]
\[
  \Leftrightarrow0=\lim_{n\rightarrow+\infty}R_{n}(x)=\lim_{n\rightarrow+\infty}\left[f(x)-\sum_{k=0}^{n}\frac{f^{(k)}(x_{0})}{k!}(x-x_{0})^{k}\right],
\]
\[
  x\in(x_{0}-R,x_{0}+R).
\]
\qedhere
\end{proof}

这个定理初看起来很简单, 但是它只是理论上的, 对一般的函数, 要判断 $\lim\limits_{n\to+\infty}R_{n}(x)=0$ 不是一件容易的事. 下面定理给出了一个能容易判断 $\lim\limits_{n\to+\infty}R_{n}(x)=0$ 的条件, 也就是给出了 $f(x)$ 能在 $(x_{0}-R,x_{0}+R)$ 上展开为 Taylor 级数的十分便于应用的充分条件.

\begin{theorem}[可 Taylor 展开的充分条件]\label{thm:14.2.5}
设 $f(x)$ 在 $(x_{0}-R,x_{0}+R)$ 上有各阶导数, 且对充分大的 $n\in\mathbb{N}$ (即 $\exists N_{0}\in\mathbb{N}$, 当 $n>N_{0}$), 均有
\[
  |f^{(n)}(x)|\leqslant M\ (\text{常数}),\quad\forall x\in(x_{0}-R,x_{0}+R),
\]
则 $f(x)$ 在 $(x_{0}-R,x_{0}+R)$ 中可 Taylor 展开.
\end{theorem}

\begin{proof}
因为 $f(x)$ 在 $(x_{0}-R,x_{0}+R)$ 上有各阶导数, 由 Taylor 公式的 Lagrange 余项可得
\[
  \begin{aligned}
    |R_{n}(x)|&=\left|\frac{f^{(n+1)}(\xi)}{(n+1)!}(x-x_{0})^{n+1}\right|\leqslant M\frac{|x-x_{0}|^{n+1}}{(n+1)!}\\
    &\leqslant M\frac{R^{n+1}}{(n+1)!}\to0\quad(n\to+\infty),\quad\forall x\in(x_{0}-R,x_{0}+R)
  \end{aligned}
\]
(由\extref{例 1.1.5}, 或用 $d'$Alembert 判别法的极限形式推得 $\sum_{n=0}^{\infty}\frac{R^{n}}{n!}$ 收敛, 从而应用定理~\ref{thm:12.1.2} 得到 $\lim\limits_{n\to+\infty}\frac{R^{n+1}}{(n+1)!}=0$). 根据夹逼定理, $\lim\limits_{n\to+\infty}R_{n}(x)=0$. 从定理~\ref{thm:14.2.4} 立知, $f(x)$ 在 $(x_{0}-R,x_{0}+R)$ 中可 Taylor 展开.
\end{proof}

\begin{example}\label{ex:14.2.3}
求下列重要初等函数的 Maclaurin 展开式，并说明收敛区间与收敛点集:

(1) $\mathrm{e}^{x}$; (2) $\sin x$; (3) $\cos x$; (4) $(1+x)^{\alpha}$, 其中 $\alpha$ 为任意实数.
\end{example}

\begin{proof}
(1) 设 $f(x)=\mathrm{e}^{x}$, 因为 $f^{(n)}(x)=(\mathrm{e}^{x})^{(n)}=\mathrm{e}^{x}$, $f^{(n)}(0)=1$, $n=0,1,2,\cdots$, 故其 Maclaurin 级数为
\[
  \mathrm{e}^{x}\sim\sum_{n=0}^{\infty}\frac{x^{n}}{n!}.
\]
又 $\forall R>0$, $|f^{(n)}(x)|=|\mathrm{e}^{x}|\leqslant\mathrm{e}^{R}$, $x\in(-R,R)$, $n=0,1,2,\cdots$. 根据定理~\ref{thm:14.2.5}, 等式
\[
  \mathrm{e}^{x}=\sum_{n=0}^{\infty}\frac{x^{n}}{n!}
\]
在 $(-R,R)$ 中成立. 由于 $R$ 是任意的, 故有
\[
  \mathrm{e}^{x}=\sum_{n=0}^{\infty}\frac{x^{n}}{n!},\quad x\in(-\infty,+\infty).
\]

(2) 设 $f(x)=\sin x$, 则 $f^{(n)}(x)=(\sin x)^{(n)}=\sin\left(x+\frac{n\pi}{2}\right)$, 于是
\[
  f^{(n)}(0)=\sin\frac{n\pi}{2}=\begin{cases}
    0, & n=2k,\\
    (-1)^{k}, & n=2k+1.
\end{cases}
\]
又因为
\[
  |f^{(n)}(0)|=|(\sin x)^{(n)}|=\left|\sin\left(x+\frac{n\pi}{2}\right)\right|\leqslant1,\quad x\in(-\infty,+\infty).
\]
根据定理~\ref{thm:14.2.5}, 有
\[
  \sin x=\sum_{k=0}^{\infty}\frac{(-1)^{k}}{(2k+1)!}x^{2k+1},\quad x\in(-\infty,+\infty).
\]

(3) 设 $f(x)=\cos x$, 则 $f^{(n)}(x)=(\cos x)^{(n)}=\cos\left(x+\frac{n\pi}{2}\right)$, 于是
\[
  f^{(n)}(0)=\cos\frac{n\pi}{2}=\begin{cases}
    (-1)^{k}, & n=2k,\\
    0, & n=2k+1.
\end{cases}
\]
又因为
\[
  |f^{(n)}(x)|=|(\cos x)^{(n)}|=\left|\cos\left(x+\frac{n\pi}{2}\right)\right|\leqslant1,\quad x\in(-\infty,+\infty).
\]
根据定理~\ref{thm:14.2.5}, 有
\[
  \cos x=\sum_{k=0}^{\infty}\frac{(-1)^{k}}{(2k)!}x^{2k},\quad x\in(-\infty,+\infty).
\]
或者由定理~\ref{thm:14.1.4}, 有
\[
  \cos x=(\sin x)'=\left(\sum_{k=0}^{\infty}\frac{(-1)^{k}}{(2k+1)!}x^{2k+1}\right)'=\sum_{k=0}^{\infty}\frac{(-1)^{k}}{(2k)!}x^{2k},\quad x\in(-\infty,+\infty).
\]

(4) 由\extref{例 4.1.6} 已经证明, $\forall x>-1$, 有
\[
  (1+x)^{\alpha}=\sum_{k=0}^{n}\mathrm{C}_{\alpha}^{k}x^{k}+R_{n}(x),
\]
其中
\[
  \mathrm{C}_{\alpha}^{k}=\frac{\alpha(\alpha-1)\cdots(\alpha-k+1)}{k!},\quad\mathrm{C}_{\alpha}^{0}=1,
\]
\[
  R_{n}(x)=\frac{\alpha(\alpha-1)\cdots(\alpha-n)}{n!}\left(\frac{1-\theta}{1+\theta x}\right)^{n}(1+\theta x)^{\alpha-1}x^{n+1},\quad0<\theta<1.
\]
对 $R_{n}(x)$ 作如下的估计:
\[
  |R_{n}(x)|\leqslant\begin{cases}
    \dfrac{|\alpha(\alpha-1)\cdots(\alpha-n)|}{n!}(1+|x|)^{\alpha-1}|x|^{n+1}, & \alpha>1,\\[2ex]
    \dfrac{|\alpha(\alpha-1)\cdots(\alpha-n)|}{n!}(1-|x|)^{\alpha-1}|x|^{n+1}, & \alpha<1.
\end{cases}
\]
根据 $d'$Alembert 判别法的极限形式知, 当 $|x|<1$ 时, 级数
\[
  \sum_{n=0}^{\infty}\frac{|\alpha(\alpha-1)\cdots(\alpha-n)|}{n!}|x|^{n+1}<+\infty.
\]
由定理~\ref{thm:12.1.2}, 有
\[
  \lim_{n\to+\infty}\frac{|\alpha(\alpha-1)\cdots(\alpha-n)|}{n!}|x|^{n+1}=0.
\]
由此及上述关于 $|R_{n}(x)|$ 的不等式推得 $\lim\limits_{n\to+\infty}R_{n}(x)=0$, $|x|<1$. 再应用定理~\ref{thm:14.2.4}, $(1+x)^{\alpha}$ 有 Taylor 展开
\[
  (1+x)^{\alpha}=\sum_{n=0}^{\infty}\mathrm{C}_{\alpha}^{n}x^{n},\quad|x|<1.
\]
右边幂级数的收敛半径
\[
  R=\frac{1}{\lim\limits_{n\to+\infty}\left|\frac{a_{n+1}}{a_{n}}\right|}=\frac{1}{\lim\limits_{n\to+\infty}\left|\frac{\mathrm{C}_{\alpha}^{n+1}}{\mathrm{C}_{\alpha}^{n}}\right|}=\frac{1}{\lim\limits_{n\to+\infty}\left|\frac{\alpha-n}{n}\right|}=1.
\]
因此, 该幂级数的收敛区间为 $(-1,1)$.

上面的 Taylor 公式在 $x=\pm1$ 处是否成立? 分下面几种情形来讨论:

① 如果 $\alpha\leqslant-1$. 此时, 有
\[
  \left|\mathrm{C}_{\alpha}^{n}(\pm1)^{n}\right|=\frac{|\alpha(\alpha-1)\cdots(\alpha-n+1)|}{n!}\geqslant\frac{n!}{n!}=1\quad(n\rightarrow+\infty),
\]
所以, 当 $x=\pm1$ 时, 幂级数
\[
  \sum_{n=0}^{\infty}\mathrm{C}_{\alpha}^{n}(\pm1)^{n}
\]
都是发散的. 因此, 该幂级数的收敛点集为 $(-1,1)$.

② 如果 $\alpha>0$. 因为
\[
  \begin{aligned}
    \frac{\mathrm{C}_{\alpha}^{n+1}(\pm1)^{n+1}}{\mathrm{C}_{\alpha}^{n}(\pm1)^{n}}&=\left|\frac{\alpha(\alpha-1)\cdots(\alpha-n)}{(n+1)!}\right|\bigg/\left|\frac{\alpha(\alpha-1)\cdots(\alpha-n+1)}{n!}\right|\\
    &=\frac{|\alpha-n|}{n+1}\to1\quad(n\to+\infty),
  \end{aligned}
\]
故用 $d'$Alembert 判别法不能断定它的敛散性. 现用 Raabe 判别法, 因为当 $n>\alpha$ 时, 有
\[
  n\left(\frac{|\mathrm{C}_{\alpha}^{n}|}{|\mathrm{C}_{\alpha}^{n+1}|}-1\right)=n\left(\frac{n+1}{n-\alpha}-1\right)=\frac{n(1+\alpha)}{n-\alpha}\to1+\alpha>1\quad(n\to+\infty),
\]
所以, $\sum_{n=0}^{\infty}|\mathrm{C}_{\alpha}^{n}(\pm1)^{n}|=\sum_{n=0}^{\infty}|\mathrm{C}_{\alpha}^{n}|<+\infty$, 收敛. 于是, 该幂级数的收敛点集为 $[-1,1]$.

③ 如果 $-1<\alpha<0$.

当 $x=-1$ 时,
\[
  \begin{aligned}
    (-1)^{n}\mathrm{C}_{\alpha}^{n}&=\frac{(-\alpha)(-\alpha+1)\cdots(-\alpha+n-1)}{n!}\\
    &=\frac{|\alpha|(|\alpha|+1)\cdots(|\alpha|+n-1)}{n!}\\
    &=\frac{|\alpha|}{n}\cdot\frac{|\alpha|+1}{1}\cdot\frac{|\alpha|+2}{2}\cdots\frac{|\alpha|+n-1}{n-1}\\
    &\geqslant\frac{|\alpha|}{n},\quad n=1,2,\cdots.
  \end{aligned}
\]
由 $\sum_{n=1}^{\infty}\frac{|\alpha|}{n}$ 发散及比较判别法知, $\sum_{n=0}^{\infty}\mathrm{C}_{\alpha}^{n}(-1)^{n}$ 发散.

当 $x=1$ 时, 级数
\[
  \sum_{n=0}^{\infty}\mathrm{C}_{\alpha}^{n}\cdot1^{n}=\sum_{n=0}^{\infty}\mathrm{C}_{\alpha}^{n}=1+\frac{\alpha(\alpha-1)}{2!}+\frac{\alpha(\alpha-1)(\alpha-2)}{3!}+\cdots
\]
为一个交错级数. 而且由于 $0<|\alpha|<1$, 故
\[
  \frac{|\alpha-n|}{n+1}\leqslant\frac{|\alpha|+n}{n+1}<1,
\]
\[
  |\mathrm{C}_{\alpha}^{n}|=\left|\frac{\alpha(\alpha-1)\cdots(\alpha-n+1)}{n!}\right|\geqslant\left|\frac{\alpha(\alpha-1)\cdots(\alpha-n+1)}{n!}\right|\cdot\left|\frac{\alpha-n}{n+1}\right|=|\mathrm{C}_{\alpha}^{n+1}|,
\]
即 $\{|\mathrm{C}_{\alpha}^{n}|\}$ 为一个单调减数列. 再根据例~\ref{ex:12.5.9} 知, $\lim\limits_{n\to+\infty}|\mathrm{C}_{\alpha}^{n}|=0$. 故由 Leibniz 判别法可断言 $\sum_{n=0}^{\infty}\mathrm{C}_{\alpha}^{n}\cdot1^{n}=\sum_{n=0}^{\infty}\mathrm{C}_{\alpha}^{n}$ 收敛. 于是, 该幂级数的收敛点集为 $(-1,1]$.

综合上述, 根据 Abel 第二定理~\ref{thm:14.1.5}, 展开式
\[
  (1+x)^{\alpha}=\sum_{n=0}^{\infty}\mathrm{C}_{\alpha}^{n}x^{n}
\]
当 $\alpha\leqslant-1$ 时, 只在 $(-1,1)$ 中成立; 当 $\alpha>0$ 时, 只在 $[-1,1]$ 中成立; 当 $-1<\alpha<0$ 时, 只在 $(-1,1]$ 中成立.

例如: 当 $\alpha=-1$ 时, 有
\[
  \begin{aligned}
    \sum_{n=0}^{\infty}(-1)^{n}x^{n}&=\sum_{n=0}^{\infty}(-x)^{n}=\frac{1}{1-(-x)}=\frac{1}{1+x}\\
    &=(1+x)^{-1}=1+\sum_{n=1}^{\infty}\frac{(-1)(-1-1)\cdots(-1-n+1)}{n!}x^{n}\\
    &=\sum_{n=0}^{\infty}(-1)^{n}x^{n},\quad x\in(-1,1).
  \end{aligned}
\]
当 $\alpha=\frac{1}{2}$ 时, 有
\[
  \begin{aligned}
    \sqrt{1+x}&=(1+x)^{\frac{1}{2}}=1+\frac{1}{2}x+\frac{\frac{1}{2}\left(-\frac{1}{2}\right)}{2!}x^{2}+\frac{\frac{1}{2}\left(-\frac{1}{2}\right)\left(-\frac{3}{2}\right)}{3!}\\
    &\quad+\cdots+\frac{\frac{1}{2}\left(-\frac{1}{2}\right)\left(-\frac{3}{2}\right)\cdots\frac{3-2n}{2}}{n!}x^{n}+\cdots\\
    &=1+\frac{x}{2}+\sum_{n=2}^{\infty}(-1)^{n-1}\frac{(2n-3)!!}{2^{n}\cdot n!}x^{n}\\
    &=1+\frac{x}{2}+\sum_{n=2}^{\infty}(-1)^{n-1}\frac{(2n-3)!!}{(2n)!!}x^{n},\quad x\in[-1,1].
  \end{aligned}
\]
当 $\alpha=-\frac{1}{2}$ 时, 有
\[
  \begin{aligned}
    \frac{1}{\sqrt{1+x}}&=(1+x)^{-\frac{1}{2}}=1+\left(-\frac{1}{2}\right)x+\frac{\left(-\frac{1}{2}\right)\left(-\frac{3}{2}\right)}{2!}x^{2}\\
    &\quad+\cdots+\frac{\left(-\frac{1}{2}\right)\left(-\frac{3}{2}\right)\cdots\left(-\frac{1}{2}-n+1\right)}{n!}x^{n}+\cdots\\
    &=1+\sum_{n=1}^{\infty}(-1)^{n}\frac{(2n-1)!!}{2^{n}\cdot n!}x^{n}=1+\sum_{n=1}^{\infty}(-1)^{n}\frac{(2n-1)!!}{(2n)!!}x^{n},\quad x\in(-1,1].
  \end{aligned}
\]
或者,
\[
  \begin{aligned}
    \frac{1}{\sqrt{1+x}}&=2\left((1+x)^{\frac{1}{2}}\right)'=2\left[1+\frac{x}{2}+\sum_{n=2}^{\infty}(-1)^{n-1}\frac{(2n-3)!!}{(2n)!!}x^{n}\right]'\\
    &=2\left[\frac{1}{2}+\sum_{n=2}^{\infty}(-1)^{n-1}\frac{(2n-3)!!}{(2n)!!}nx^{n-1}\right]\\
    &=1+\sum_{n=2}^{\infty}(-1)^{n-1}\frac{(2n-3)!!}{(2n-2)!!}x^{n-1}\\
    &=1+\sum_{n=1}^{\infty}(-1)^{n}\frac{(2n-1)!!}{(2n)!!}x^{n},\quad x\in(-1,1].
  \end{aligned}
\]
\end{proof}

提醒读者, 例~\ref{ex:14.1.7} 及例~\ref{ex:14.2.3} 中 $6$ 个初等函数的 Maclaurin 展开式必须牢牢地熟记!

定理~\ref{thm:14.1.9} 与定理~\ref{thm:14.2.5} 给出了将函数展开为幂级数的方法. 除此之外, 利用某些函数的已知展开式, 通过幂级数的求导, 积分以及代数运算也能作出其他一些函数的幂级数展开式.

例~\ref{ex:14.1.7} 通过幂级数的逐项积分求得了函数 $\ln(1+x)$ 与 $\arctan x$ 的幂级数展开式. 读者回顾有限幂级数展开:
\[
  f(x)=\sum_{k=0}^{n}a_{k}(x-x_{0})^{k}+o((x-x_{0})^{n}).
\]
根据惟一性定理\extref{定理 4.1.1} 以及\extref{定理 4.1.2}, 它们表明无论用什么方法, 所得到的有限幂级数展开一定是有限 Taylor 展开:
\[
  f(x)=\sum_{k=0}^{n}\frac{f^{(k)}(x_{0})}{k!}(x-x_{0})^{k}+o((x-x_{0})^{n}).
\]
因此, 只要有什么展开的妙法, 尽管放心使用! \extref{例 4.1.7}$\sim$\extref{例 4.1.13} 就体现了这一思想.

进而, 定理~\ref{thm:14.2.1} 说明收敛幂级数展开也是惟一的, 它就是收敛的 (无穷) Taylor 展开:
\[
  f(x)=\sum_{n=0}^{\infty}\frac{f^{(n)}(x)}{n!}(x-x_{0})^{n}.
\]
因此, 无论用什么方法, 其最后的展开是一样的, 都是 (无穷) Taylor 展开.

%  第二部分：例 14.2.4 起（生成脚本 tools/_make_ch14_2b.py）
% ==================================================================

% ===== PART-B =====

\begin{example}\label{ex:14.2.4}
将 $\dfrac{1}{(1-x)(2-x)}$ 展开为 Maclaurin 级数.
\end{example}

\[
  \begin{aligned}
    \frac{1}{(1-x)(2-x)}&=\frac{1}{1-x}-\frac{1}{2-x}=\frac{1}{1-x}-\frac{1}{2}\frac{1}{1-\frac{x}{2}}\\
    &=\sum_{n=0}^{\infty}x^{n}-\frac{1}{2}\sum_{n=0}^{\infty}\left(\frac{x}{2}\right)^{n}=\sum_{n=0}^{\infty}\left(1-\frac{1}{2^{n+1}}\right)x^{n},\quad x\in(-1,1),
  \end{aligned}
\]
其中第 $1$ 个幂级数的收敛区间为 $(-1,1)$, 第 $2$ 个幂级数的收敛区间为 $(-2,2)$.

因为
\[
  \left|\left(1-\frac{1}{2^{n+1}}\right)(\pm1)^{n}\right|=1-\frac{1}{2^{n+1}}\to1-0=1\quad(n\to+\infty),
\]
所以
\[
  \sum_{n=0}^{\infty}\left(1-\frac{1}{2^{n+1}}\right)(\pm1)^{n}
\]
都发散. 由此推得上述幂级数的收敛点集为 $(-1,1)$.

\begin{example}\label{ex:14.2.5}
求 (1) $\mathrm{e}^{-x^{2}}$; (2) $\sin x^{3}$ 的 Maclaurin 展开.
\end{example}

\begin{proof}
(1) $\mathrm{e}^{-x^{2}}=\sum_{n=0}^{\infty}\frac{(-x^{2})^{n}}{n!}=\sum_{n=0}^{\infty}\frac{(-1)^{n}x^{2n}}{n!}$, $x\in(-\infty,+\infty)$.

(2) $\sin x^{3}=\sum_{k=0}^{\infty}\frac{(-1)^{k}}{(2k+1)!}(x^{3})^{2k+1}=\sum_{k=0}^{\infty}\frac{(-1)^{k}}{(2k+1)!}x^{6k+3}$, $x\in(-\infty,+\infty)$.
\end{proof}

\begin{example}\label{ex:14.2.6}
设幂级数 $\sum_{n=0}^{\infty}a_{n}x^{n}$ 的收敛半径为 $1$, 则:

(1) $\dfrac{1}{1-x}\sum_{n=0}^{\infty}a_{n}x^{n}=\sum_{n=0}^{\infty}S_{n}x^{n}$, $x\in(-1,1)$, 其中 $S_{n}=\sum_{k=0}^{n}a_{k}$;
\[
  \frac{1}{(1-x)^{2}}\sum_{n=0}^{\infty}a_{n}x^{n}=\sum_{n=0}^{\infty}(S_{0}+S_{1}+\cdots+S_{n})x^{n},\quad x\in(-1,1).
\]
\end{example}

\begin{proof}
(1)
\[
  \begin{aligned}
    \frac{1}{1-x}\sum_{n=0}^{\infty}a_{n}x^{n}&=(1+x+x^{2}+\cdots)(a_{0}+a_{1}x+a_{2}x^{2}+\cdots)\\
    &=a_{0}+(a_{0}+a_{1})x+(a_{0}+a_{1}+a_{2})x^{2}\\
    &\quad+\cdots+(a_{0}+a_{1}+\cdots+a_{n})x^{n}+\cdots\\
    &=\sum_{n=0}^{\infty}S_{n}x^{n}.
  \end{aligned}
\]
或者令 $b_{m}=1$, $m=0,1,2,\cdots$, 则由定理~\ref{thm:14.1.8}(2),
\[
  \begin{aligned}
    \frac{1}{1-x}\sum_{n=0}^{\infty}a_{n}x^{n}&=\left(\sum_{m=0}^{\infty}b_{m}x^{m}\right)\left(\sum_{k=0}^{\infty}a_{k}x^{k}\right)\\
    &=\sum_{n=0}^{\infty}\left(\sum_{k=0}^{n}b_{n-k}a_{k}\right)x^{n}=\sum_{n=0}^{\infty}\left(\sum_{k=0}^{n}a_{k}\right)x^{n}\\
    &=\sum_{n=0}^{\infty}S_{n}x^{n}.
  \end{aligned}
\]

(2)
\[
  \begin{aligned}
    \frac{1}{(1-x)^{2}}\sum_{n=0}^{\infty}a_{n}x^{n}&=\frac{1}{1-x}\left(\frac{1}{1-x}\sum_{n=0}^{\infty}a_{n}x^{n}\right)\xlongequal{\text{由 (1)}}\frac{1}{1-x}\sum_{n=0}^{\infty}S_{n}x^{n}\\
    &\xlongequal{\text{由 (1)}}\sum_{n=0}^{\infty}(S_{0}+\cdots+S_{n})x^{n}.
  \end{aligned}
\]
\end{proof}

\begin{example}\label{ex:14.2.7}
将函数 $\dfrac{\ln(1-x)}{1-x}$ 展开为 Maclaurin 级数.
\end{example}

\begin{proof}
由例~\ref{ex:14.1.7}, 有
\[
  \ln(1-x)=\sum_{n=1}^{\infty}\frac{(-1)^{n-1}}{n}(-x)^{n}=-\sum_{n=1}^{\infty}\frac{x^{n}}{n},\quad-1\leqslant x<1.
\]
根据例~\ref{ex:14.2.6} 立即可得
\[
  \frac{\ln(1-x)}{1-x}=\frac{1}{1-x}\left(-\sum_{n=1}^{\infty}\frac{x^{n}}{n}\right)=-\sum_{n=1}^{\infty}\left(1+\frac{1}{2}+\cdots+\frac{1}{n}\right)x^{n},\quad-1<x<1.
\]
\end{proof}

\begin{example}\label{ex:14.2.8}
设 $f(x)$ 及其所有导数在 $[0,r]$ 非负, 则 $f(x)$ 的 Taylor 级数在 $[0,r]$ 上收敛于 $f(x)$, 即
\[
  \sum_{n=0}^{\infty}\frac{f^{(n)}(0)}{n!}x^{n}=f(x)\quad(0\leqslant x<r).
\]
\end{example}

\begin{proof}
令 $R_{n}(x)=f(x)-\sum_{k=0}^{n}\frac{f^{(k)}(0)}{k!}x^{k}$, $g(x)=\dfrac{R_{n}(x)}{x^{n+1}}$, 则
\[
  g'(x)=\frac{xR_{n}'(x)-(n+1)R_{n}(x)}{x^{n+2}}.
\]
再令 $h(x)=xR_{n}'(x)-(n+1)R_{n}(x)$, 逐次求导得到
\[
  h'(x)=xR_{n}''(x)-nR_{n}'(x),
\]
\[
  \bullet\quad\bullet
\]
\[
  h^{(n)}(x)=xR_{n}^{(n+1)}(x)-R_{n}^{(n)}(x),
\]
\[
  h^{(n+1)}(x)=xR_{n}^{(n+2)}(x)=xf^{(n+2)}(x)\geqslant0.
\]
于是, $h^{(n)}(x)\geqslant h^{(n)}(0)=0$, 又得 $h^{(n-1)}(x)\geqslant h^{(n-1)}(0)=0$, $\cdots$, 继续推出 $h(x)\geqslant0$, 进而有 $g'(x)\geqslant0$, 所以 $g(x)\leqslant g(r)$, 且
\[
  0\leqslant\frac{f^{(n+1)}(\xi)}{(n+1)!}x^{n+1}=R_{n}(x)\leqslant\left(\frac{x}{r}\right)^{n+1}f(r)\to0\quad(n\to+\infty).
\]
因此, $\lim\limits_{n\to+\infty}R_{n}(x)=0$, 于是
\[
  f(x)=\sum_{n=0}^{\infty}\frac{f^{(n)}(0)}{n!}x^{n},\quad0\leqslant x<r.
\]
\end{proof}

\begin{exblock}{练习题 14.2}
\begin{enumerate}
  \item 利用已知的初等函数展开式，给出下列函数的幂级数展开式:
  \begin{enumerate}
    \item $\mathrm{e}^{x^{2}}$;
    \item $\cos^{2}x$;
    \item $\dfrac{x^{12}}{1-x}$;
    \item $\ln\sqrt{\dfrac{1+x}{1-x}}$;
    \item $\dfrac{x}{1+x-2x^{2}}$;
    \item $(1+x)\mathrm{e}^{-x}$.
  \end{enumerate}
  \item 求下列函数的幂级数展开式:
  \begin{enumerate}
    \item $(1+x)\ln(1+x)$;
    \item $\arctan x-\ln\sqrt{1+x^{2}}$;
    \item $\dfrac{x}{(1-x)(1-x^{2})}$;
    \item $(1+x^{2})\arctan x$.
  \end{enumerate}
  \item 将下列函数展开成幂级数:
  \begin{enumerate}
    \item $\arcsin x$;
    \item $\ln(x+\sqrt{1+x^{2}})$;
    \item $\displaystyle\int_{0}^{x}\frac{\sin t}{t}\mathrm{d}t$.
  \end{enumerate}
  \item 将函数 $f(x)=\dfrac{x}{\sqrt{1+x}}$ 按 $\dfrac{x}{1+x}$ 的正整数幂展开成幂级数.
  \item 将下列级数展开为 $x$ 的幂级数:

  (1) $f(x)=\dfrac{1}{(1+x)(1+x^{2})(1+x^{4})}$ (提示: $f(x)=\dfrac{1-x}{(1-x)(1+x)(1+x^{2})(1+x^{4})}$);

  (2) $f(x)=\sin^{3}x$ (提示: $f(x)=\dfrac{3}{4}\sin x-\dfrac{1}{4}\sin3x$).
  \item 设 $x>0$, 证明:
  \[
    \ln x=2\left[\frac{x-1}{x+1}+\frac{1}{3}\left(\frac{x-1}{x+1}\right)^{3}+\frac{1}{5}\left(\frac{x-1}{x+1}\right)^{5}+\cdots\right]
\]
  (提示: 令 $t=\dfrac{x-1}{x+1}$, 即 $x=\dfrac{1+t}{1-t}$).
  \item 应用逐项积分证明:

  (1) 函数 $f(x)=\arctan\dfrac{2x}{2-x^{2}}$ 的幂级数展开式为
  \[
    f(x)=\int_{0}^{x}f'(t)\mathrm{d}t=\sum_{n=0}^{\infty}(-1)^{[\frac{n}{2}]}\frac{x^{2n+1}}{2^{n}(2n+1)},\quad|x|<\sqrt{2};
\]
  (2) 上面展开式的收敛点集为 $[-\sqrt{2},\sqrt{2}]$.
  \item 将 $f(x)=\dfrac{1}{\mathrm{e}}(1+x)^{\frac{1}{x}}$ 按 $x$ 的幂展开至三次项 (提示: $f(x)=\mathrm{e}^{\frac{1}{x}\ln(1+x)-1}$, 并应用 $\ln(1+x)$ 与 $\mathrm{e}^{x}$ 的展开式).
\end{enumerate}
\end{exblock}

\begin{exblock}{思考题 14.2}
\begin{enumerate}
  \item 证明: $f(x)=\ln^{2}(1-x)$ 的幂级数展开式为
  \[
    f(x)=2\sum_{n=1}^{\infty}\left(1+\frac{1}{2}+\frac{1}{3}+\cdots+\frac{1}{n}\right)\frac{x^{n+1}}{n+1},\quad x\in[-1,1)
\]
  (提示: (1) 应用 $f(x)=\ln(1-x)\cdot\ln(1-x)$ 的幂级数乘法; (2) 先求 $f'(x)=-\dfrac{2\ln(1-x)}{1-x}$ 的展开式).
  \item 按下面两种方法将
  \[
    y=y(x)=\frac{\ln(x+\sqrt{1+x^{2}})}{\sqrt{1+x^{2}}}
\]
  展开为幂级数
  \[
    y=\frac{\ln(x+\sqrt{1+x^{2}})}{\sqrt{1+x^{2}}}=\sum_{n=0}^{\infty}(-1)^{n}\frac{(2n)!!}{(2n+1)!!}x^{2n+1},\quad x\in[-1,1].
\]
  (1) 先证: $(1+x^{2})y^{(n+1)}+(2n+1)xy^{(n)}+n^{2}y^{(n-1)}=0$,
  \[
    y^{(2n)}(0)=0,\ n=1,2,\cdots,
\]
  \[
    y^{(2n+1)}(0)=(-1)^{n}[(2n)!!]^{2};
\]
  (2) 由 $(1+x^{2})y'=1-xy$, $y=\sum_{n=0}^{\infty}a_{n}x^{n}$ 以及待定系数法证明
  \[
    a_{2n}=0,\quad a_{2n+1}=(-1)^{n}\frac{(2n)!!}{(2n+1)!!}.
\]
  \item (1) 对 $\sum_{n=0}^{\infty}(-1)^{n}x^{n}=\dfrac{1}{1+x}$ ($|x|<1$) 两边从 $0$ 到 $x$ 两次积分, 证明:
  \[
    \sum_{n=0}^{\infty}(-1)^{n}\frac{x^{n+2}}{(n+1)(n+2)}=x\ln(1+x)-x+\ln(1+x),\quad|x|<1;
\]
  (2) 证明: $\displaystyle\sum_{n=0}^{\infty}\frac{(-1)^{n}}{(n+1)(n+2)}=2\ln2-1$.
  \item 设 $k$ 为任一正整数, 证明:

  (1) $f(x)=\displaystyle\sum_{n=1}^{\infty}\frac{n^{k}}{n!}x^{n}$ 的收敛半径 $R=+\infty$;

  (2) $\displaystyle\int_{0}^{x}\frac{1}{t}f(t)\mathrm{d}t=\sum_{n=1}^{\infty}\frac{n^{k-1}}{n!}x^{n}$;

  (3) $\displaystyle\int_{0}^{x}\frac{1}{t_{k-1}}\mathrm{d}t_{k-1}\int_{0}^{t_{k-1}}\frac{1}{t_{k-2}}\mathrm{d}t_{k-2}\cdots\int_{0}^{t_{1}}\frac{1}{t}f(t)\mathrm{d}t=\mathrm{e}^{x}-1$;

  (4) $f(x)=p_{k}(x)\mathrm{e}^{x}$, 其中 $p_{k}(x)$ 为 $x$ 的整系数 $k$ 次多项式;

  (5) $\displaystyle\sum_{n=1}^{\infty}\frac{n^{k}}{n!}$ 为 $\mathrm{e}$ 的整数倍.
  \item 设幂级数为
  \[
    S(x)=1+x+\frac{x^{2}}{2}+\frac{x^{3}}{1\times3}+\frac{x^{4}}{2\times4}+\frac{x^{5}}{1\times3\times5}+\frac{x^{6}}{2\times4\times6}+\cdots,
\]
  证明:

  (1) 收敛半径 $R=+\infty$;

  (2) $S'(x)=1+xS(x)$;

  (3) $S(x)=\mathrm{e}^{\frac{x^{2}}{2}}\left(\int_{0}^{x}\mathrm{e}^{-\frac{t^{2}}{2}}\mathrm{d}t+1\right)$.
  \item 应用幂级数展开证明: 不等式
  \[
    \mathrm{e}^{x}+\mathrm{e}^{-x}\leqslant2\mathrm{e}^{\frac{x^{2}}{2}},\quad x\in(-\infty,+\infty).
\]
\end{enumerate}
\end{exblock}
